\SourceLeaf{322}
Quod igitur in exemplis artt.\ 273, 274 evenit, ut forma definita
$\left(\begin{smallmatrix}19,21,50\\15,28,1\end{smallmatrix}\right)$
determinantis $-1$ ad hanc $xx+yy+zz$, atque forma indefinita
$\left(\begin{smallmatrix}10,26,2\\7,0,4\end{smallmatrix}\right)$
determinantis $2$ ad $2xx-2yz$ sive (quod eodem redit) ad $2xx+2yz$ reduceretur,
per disquisitiones praecedentes a priori praevideri potuisset.

\Needspace{8\baselineskip}
\Article{278}
Per formam ternariam, cuius indeterminatae sunt $x,x',x''$, \textit{repraesentantur}
tum numeri, tribuendo ipsis $x,x',x''$ valores determinatos, tum formae binariae
per huiusmodi substitutiones
\[
x=mt+nu,\qquad x'=m't+n'u,\qquad x''=m''t+n''u
\]
designantibus $m,n,m'$ etc.\ numeros determinatos; $t,u$ indeterminatas formae
repraesentatae. Ad theoriam itaque completam formarum ternariarum requireretur
solutio sequentium problematum: I.\ Invenire omnes repraesentationes numeri dati
per formam ternariam datam. II.\ Invenire omnes repraesentationes formae binariae
datae per ternariam datam. III.\ Diiudicare, utrum duae formae ternariae datae
eiusdem determinantis aequivalentes sint necne, et in casu priori omnes
transformationes alterius in alteram invenire. IV.\ Diiudicare, utrum forma
ternaria data aliam datam determinantis maioris implicet, necne, et in casu
priori omnes transformationes illius in hanc assignare. De quibus problematibus
longe difficilioribus quam analoga in formis binariis alio loco pluribus agemus:
hic disquisitionem nostram restringimus ad ostendendum, quomodo problema primum
ad secundum secundumque ad tertium reduci possit; tertium vero pro casibus
quibusdam simplicissimis formarumque binariarum theoriam imprimis illustrantibus
solvere docebimus; quartum hic omnino excludemus.

\Needspace{8\baselineskip}
\Article{279}
\textsc{Lemma.} \textit{Propositis tribus numeris integris quibuscunque $a,a',a''$
(qui tamen non omnes simul $=0$): invenire sex alios $B,B',B'',C,C',C''$ ita
comparatos ut fiat}
\[
B'C''-B''C'=a,\qquad B''C-BC''=a',\qquad BC'-B'C=a''
\]

\textit{Sol.} Sit $\alpha$ div.\ comm.\ max.\ ipsorum $a,a',a''$. accipianturque
integri $A,A',A''$ ita ut fiat
\[
Aa+A'a'+A''a''=\alpha
\]
\SourceLeaf{323}
Porro accipiantur tres integri $\mathfrak{C},\allowbreak \mathfrak{C}',\allowbreak \mathfrak{C}''$ ad
lubitum ea sola conditione, ut tres numeri
$\mathfrak{C}'A''-\mathfrak{C}''A',\allowbreak \ \mathfrak{C}''A-\mathfrak{C}A'',\allowbreak \
\mathfrak{C}A'-\mathfrak{C}'A$, quos resp.\ per $b,b',b''$ ipsorumque divisorem
communem maximum per $\beta$ designabimus, non fiant simul $=0$. Tunc ponatur
\[
a'b''-a''b'=\alpha\beta C,\qquad
a''b-ab''=\alpha\beta C',\qquad
ab'-a'b=\alpha\beta C''
\]
patetque, ipsos $C,C',C''$ fore integros. Denique accipiendo integros
$\mathfrak{B},\allowbreak \mathfrak{B}',\allowbreak \mathfrak{B}''$ ita ut fiat
\[
\mathfrak{B}b+\mathfrak{B}'b'+\mathfrak{B}''b''=\beta
\]
ponendo
\[
\mathfrak{B}a+\mathfrak{B}'a'+\mathfrak{B}''a''=h
\]
et statuendo
\[
B=\alpha\mathfrak{B}-hA,\qquad
B'=\alpha\mathfrak{B}'-hA',\qquad
B''=\alpha\mathfrak{B}''-hA''
\]
hi valores ipsorum $B,B',B'',C,C',C''$ aequationibus praescriptis satisfacient.

Invenitur enim
\[
\begin{gathered}
aB+a'B'+a''B''=0\\
bA+b'A'+b''A''=0\quad\text{unde}\quad bB+b'B'+b''B''=\alpha\beta
\end{gathered}
\]
Iam ex valoribus ipsorum $C',C''$ fit
\[
\begin{aligned}
\alpha\beta(B'C''-B''C')
&=ab'B'-a'bB'-a''bB''+ab''B''\\
&=a(bB+b'B'+b''B'')-b(aB+a'B'+a''B'')=\alpha\beta a
\end{aligned}
\]
adeoque $B'C''-B''C'=a$; similique modo invenitur $B''C-BC''=a'$, $BC'-B'C=a''$.
\textit{Q. E. F.} --- Ceterum analysis per quam haec solutio inventa est,
nec non methodus ex una solutione omnes inveniendi, hic sunt supprimendae.

\Needspace{8\baselineskip}
\Article{280}
Supponamus, formam binariam
\[
att+2btu+cuu\ldots\varphi
\]
\SourceLeaf{324}
cuius determinans $=D$, repraesentari per formam ternariam $f$, cuius
indeterminatae $x,x',x''$, ponendo
\[
x=mt+nu,\qquad x'=m't+n'u,\qquad x''=m''t+n''u
\]
ipsique $f$ adiunctam esse formam $F$, cuius indeterminatae $X,X',X''$. Tunc
per calculum facile confirmatur (designando coëfficientes formarum $f,F$ per
literas peculiares) sive etiam ex art.\ 268.\ II.\ protinus deducitur, numerum
$D$ repraesentari per $F$ ponendo
\[
X=m'n''-m''n',\qquad X'=m''n-mn'',\qquad X''=mn'-m'n
\]
quae repraesentatio numeri $D$ repraesentationi formae $\varphi$ per $f$
\textit{adiuncta} commode dici potest. Si valores ipsarum $X,X',X''$ divisorem
communem non habent, brevitatis caussa hanc repraesentationem ipsius $D$
\textit{propriam} vocabimus, sin secus, \textit{impropriam}, easdem denominationes
etiam repraesentationi formae $\varphi$ per $f$, cui illa repraes.\ ipsius $D$
adiuncta est, tribuemus. Iam inventio omnium repraesentationum propriarum numeri
$D$ per formam $F$ sequentibus momentis innititur:

I.\ Nulla repraesentatio ipsius $D$ per $F$ datur, quae non ex aliqua
repraesentatione alicuius formae determinantis $D$ per formam $f$ deduci possit,
\textit{i. e.}\ tali repraesentationi adiuncta sit.

Sit enim repraesentatio quaecunque ipsius $D$ per $F$ haec: $X=L$, $X'=L'$,
$X''=L''$; accipiantur per lemma art.\ praec.\ $m,m',m'',n,n',n''$ ita ut fiat
\[
m'n''-m''n'=L,\qquad m''n-mn''=L',\qquad mn'-m'n=L''
\]
transeatque $f$ per substitutionem
\[
x=mt+nu,\qquad x'=m't+n'u,\qquad x''=m''t+n''u
\]
in formam binariam $\varphi=att+2btu+cuu$. Tunc facile perspicietur, $D$ fore
determinantem formae $\varphi$ ipsiusque repraesentationi per $f$ repraesentationem
propositam ipsius $D$ per $F$ adiunctam.

\textit{Ex.} Sit $f=xx+x'x'+x''x''$, adeoque $F=-XX-X'X'-X''X''$;
$D=-209$; ipsiusque repraesentatio per $F$ haec $X=1$, $X'=8$, $X''=12$;
\SourceLeaf{325}
hinc inveniuntur valores ipsorum $m,m',m'',n,n',n''$ hi $-20,1,1,-12,0,1$
resp., atque $\varphi=402tt+482tu+145uu$.

II.\ Si $\varphi,\chi$ sunt formae binariae proprie aequivalentes, quaevis
repraesentatio ipsius $D$ per $F$ alicui repraesentationi formae $\varphi$ per $f$
adiuncta, etiam alicui repraesentationi formae $\chi$ per $f$ adiuncta erit.

Sint $p,q$ indeterminatae formae $\chi$; transeat $\varphi$ in $\chi$ per
substitutionem propriam $t=\alpha p+\beta q$, $u=\gamma p+\delta q$, sitque
aliqua repraesentatio formae $\varphi$ per $f$ haec
\[
x=mt+nu,\qquad x'=m't+n'u,\qquad x''=m''t+n''u\ldots(R)
\]
Tunc nullo negotio perspicitur, si ponatur
\[
\begin{aligned}
\alpha m+\gamma n&=g,&\quad \alpha m'+\gamma n'&=g',&\quad \alpha m''+\gamma n''&=g''\\
\beta m+\delta n&=h,&\quad \beta m'+\delta n'&=h',&\quad \beta m''+\delta n''&=h''
\end{aligned}
\]
formam $\chi$ repraesentatum iri per $f$ statuendo
\[
x=gp+hq,\qquad x'=g'p+h'q,\qquad x''=g''p+h''q\ldots(R')
\]
calculoque facto invenitur (propter $\alpha\delta-\beta\gamma=1$) esse
\[
g'h''-g''h'=m'n''-m''n',\quad
g''h-gh''=m''n-mn'',\quad
gh'-g'h=mn'-m'n
\]
\textit{i. e.}\ repraesentationibus $R,R'$ eadem repraesentatio ipsius $D$ per
$F$ adiuncta est.

Ita in ex.\ praec.\ formae $\varphi$ aequivalere invenitur
$\chi=13pp-10pq+18qq$, in quam illa transit per substitutionem propriam
$t=-3p+q$, $u=5p-2q$; hinc invenitur repraesentatio formae $\chi$ per $f$ haec
$x=4q$, $x'=-3p+q$, $x''=2p-q$, ex qua eadem numeri $-209$ repraesentatio
deducitur, a qua profecti eramus.

III.\ Denique si duae formae binariae $\varphi,\chi$ determinantis $D$, quarum
indeterminatae sunt $t,u$; $p,q$, per $f$ repraesentari possunt, alicuique
repraesentationi unius eadem repraesentatio propria ipsius $D$ per $F$ adiuncta
est, atque alicui repraesentationi alterius, illae formae necessario erunt
proprie aequivalentes. Supponamus $\varphi$ repraesentari per $f$ ponendo
\[
x=mt+nu,\qquad x'=m't+n'u,\qquad x''=m''t+n''u
\]
\SourceLeaf{326}
$\chi$ vero statuendo
\[
x=gp+hq.\qquad x'=g'p+h'q,\qquad x''=g''p+h''q
\]
atque esse
\[
\begin{aligned}
m'n''-m''n'&=g'h''-g''h'=L\\
m''n-mn''&=g''h-gh''=L'\\
mn'-m'n&=gh'-g'h=L''
\end{aligned}
\]
Accipiantur integri $l,l',l''$ ita ut fiat $Ll+L'l'+L''l''=1$, ponaturque
\[
\begin{aligned}
n'l''-n''l'&=M,&\quad n''l-nl''&=M',&\quad nl'-n'l&=M''\\
l'm''-l''m'&=N,&\quad l''m-lm''&=N',&\quad lm'-l'm&=N''
\end{aligned}
\]
denique statuatur
\[
\begin{aligned}
gM+g'M'+g''M''&=\alpha,&\quad hM+h'M'+h''M''&=\beta\\
gN+g'N'+g''N''&=\gamma,&\quad hN+h'N'+h''N''&=\delta
\end{aligned}
\]
Hinc facile deducitur
\[
\begin{aligned}
\alpha m+\gamma n&=g-l(gL+g'L'+g''L'')=g\\
\beta m+\delta n&=h-l(hL+h'L'+h''L'')=h
\end{aligned}
\]
similique modo
\[
\alpha m'+\gamma n'=g',\quad \beta m'+\delta n'=h',\quad
\alpha m''+\gamma n''=g'',\quad \beta m''+\delta n''=h''
\]
Hinc patet. $mt+nu$, $m't+n'u$, $m''t+n''u$ transire per substitutionem
\[
t=\alpha p+\beta q,\quad u=\gamma p+\delta q\ldots(S)
\]
in $gp+hq$, $g'p+h'q$, $g''p+h''q$ resp., unde manifestum est, $\varphi$
transire per substitutionem $S$ in eandem formam, in quam $f$ transeat ponendo
\[
x=gp+hq,\qquad x'=g'p+h'q,\qquad x''=g''p+h''q
\]
adeoque in formam $\chi$, cui itaque aequivalet. Denique per substitutiones
debitas facile invenitur
\[
\alpha\delta-\beta\gamma=(Ll+L'l'+L''l'')^2=1
\]
quocirca substitutio $S$ est propria, formaeque $\varphi,\chi$ proprie aequivalentes.
\SourceLeaf{327}

Ex his observationibus derivantur regulae sequentes ad inveniendum omnes
repraesentationes proprias ipsius $D$ per $F$: Evolvantur omnes classes formarum
binariarum determinantis $D$, et ex singulis una forma ad libitum eligatur;
quaerantur omnes repraesentationes propriae singularum harum formarum per $f$
(reiectis iis, quae forte per $f$ repraesentari nequeunt), et ex singulis hisce
repraesentationibus deducantur repraesentationes numeri $D$ per $F$. Ex I et II
manifestum est, hoc modo omnes repraesentationes proprias possibiles obtineri,
adeoque solutionem esse completam; ex III, transformationes formarum e classibus
diversis certo producere repraesentationes diversas.

\Needspace{8\baselineskip}
\Article{281}
Investigatio repraesentationum \textit{impropriarum} numeri dati $D$ per formam
$F$ ad casum praecedentem facile reducitur. Scilicet manifestum est, si $D$ per
nullum quadratum (praeter $1$) divisibilis sit, tales repraesentationes omnino non
dari; sin secus, metientibus ipsum $D$ quadratis $\lambda\lambda,\mu\mu,\nu\nu$
etc., omnes repraesentationes improprias ipsius $D$ per $F$ inveniri, si omnes
repraesentationes propriae numerorum $\frac{D}{\lambda\lambda}$,
$\frac{D}{\mu\mu}$, $\frac{D}{\nu\nu}$ etc.\ per eandem formam evolvantur,
indeterminatarumque valores per $\lambda,\mu,\nu$ etc.\ resp.\ multiplicentur.

Hoc itaque modo inventio omnium repraesentationum numeri dati per formam ternariam
datam, \textit{quae alicui formae ternariae adiuncta est}, a problemate secundo
pendet; ad hunc vero casum, qui primo aspectu minus late patere videri posset,
reliqui ita reducuntur. Sit $D$ numerus repraesentandus per formam
$\left(\begin{smallmatrix}g,g',g''\\h,h',h''\end{smallmatrix}\right)$,
cuius determinans $\Delta$, et cui adiuncta est forma
$\left(\begin{smallmatrix}G,G',G''\\H,H',H''\end{smallmatrix}\right)=f$.
Tunc huic rursus adiuncta erit
$\left(\begin{smallmatrix}\Delta g,\Delta g',\Delta g''\\\Delta h,\Delta h',\Delta h''\end{smallmatrix}\right)=F$,
patetque, repraesentationes numeri $\Delta D$ per $F$ (quarum investigatio a
praecc.\ pendet) omnino identicas esse cum repraesentationibus numeri $D$ per
formam propositam. --- Ceterum quando omnes coëfficientes formae $f$ divisorem
communem $\mu$ habent, perspicuum est, omnes coëfficientes formae $F$ divisibiles
esse per $\mu\mu$, quocirca etiam $\Delta D$ per $\mu\mu$ divisibilis esse
debebit (alioquin nullae repraesentationes darentur); repraesentationesque numeri
$D$ per formam propositam coincident cum repraesentationibus numeri
$\frac{\Delta D}{\mu\mu}$ per formam, quae oritur ex $F$, dividendo singulos
coëfficientes per $\mu\mu$, quae forma adiuncta erit ei, quae oritur ex $f$,
dividendo singulos coëfficientes per $\mu$.

Denique observamus, hanc problematis primi solutionem in unico casu, ubi
\SourceLeaf{328}
$D=0$, non esse applicabilem; hic enim omnes formae binariae determinantis $D$
in multitudinem finitam classium non distribuuntur; infra autem hunc casum ex
aliis principiis solvemus.

\Needspace{8\baselineskip}
\Article{282}
Investigatio repraesentationum formae binariae datae, cuius determinans non $=0$\HistoricalNote{*)}{Hunc casum per methodum aliquantum diversam tractandum hoc loco brevitatis caussa praeterimus.},
per ternariam datam pendet ab observationibus sequentibus:

I.\ Ex quavis repraesentatione propria formae binariae $(p,q,r)=\varphi$
determinantis $D$ per ternariam $f$ determinantis $\Delta$ deduci possunt integri
$B,B'$ tales ut sit
\[
BB\equiv\Delta p,\qquad BB'\equiv-\Delta q,\qquad
B'B'\equiv\Delta r\ (\mathrm{mod.}\ D)
\]
\textit{i. e.}\ valor expressionis $\surd\Delta(p,\allowbreak -q,\allowbreak r)(\mathrm{mod.}\ D)$.
Habeatur repraesentatio propria formae $\varphi$ per $f$ haec
\[
x=\alpha t+\beta u,\qquad x'=\alpha't+\beta'u,\qquad x''=\alpha''t+\beta''u
\]
(designantibus $x,x',x''$; $t,u$ indeterminatas formarum $f,\varphi$); accipiantur
integri $\gamma,\gamma',\gamma''$ ita ut
\[
(\alpha'\beta''-\alpha''\beta')\gamma+
(\alpha''\beta-\alpha\beta'')\gamma'+
(\alpha\beta'-\alpha'\beta)\gamma''
\]
$=k$ fiat vel $=+1$ vel $=-1$, transeatque $f$ per substitutionem
\[
\begin{array}{ccc}
\alpha,&\beta,&\gamma\\
\alpha',&\beta',&\gamma'\\
\alpha'',&\beta'',&\gamma''
\end{array}
\]
in formam
$\left(\begin{smallmatrix}a,a',a''\\b,b',b''\end{smallmatrix}\right)=g$,
cui adiuncta sit
$\left(\begin{smallmatrix}A,A',A''\\B,B',B''\end{smallmatrix}\right)=G$.
Tunc manifestum est, fore $a=p$, $b''=q$, $a'=r$, $A''=D$, atque $\Delta$
determinantem formae $g$; unde
\[
BB=\Delta p+A'D,\qquad BB'=-\Delta q+B''D,\qquad B'B'=\Delta r+AD
\]
Ita \textit{e. g.}\ forma $19tt+6tu+41uu$ repraesentatur per $xx+x'x'+x''x''$
ponendo $x=3t+5u$, $x'=3t-4u$, $x''=t$; unde statuendo $\gamma=-1$,
$\gamma'=1$, $\gamma''=0$,
\SourceLeaf{329}
eruitur $B=-171$, $B'=27$, sive valor $(-171,27)$ expr.\
$\surd-1(19,\allowbreak -3,\allowbreak 41)(\mathrm{mod.}\ 770)$.

Hinc iam sequitur, si $\Delta(p,-q,r)$ non sit residuum quadraticum ipsius $D$,
$\varphi$ per nullam formam ternariam determinantis $\Delta$ proprie
repraesentabilem esse posse; in eo itaque casu, ubi $\Delta,D$ inter se primi
sunt, $\Delta$ numerus characteristicus formae $\varphi$ esse debebit.

II.\ Quum $\gamma,\gamma',\gamma''$ infinite multis modis diversis determinari
possint, etiam alii atque alii valores ipsorum $B,B'$ inde prodibunt, qui quem
nexum inter se habeant videamus. Ponamus, etiam $\delta,\delta',\delta''$ ita
comparatos esse, ut
\[
(\alpha'\beta''-\alpha''\beta')\delta+
(\alpha''\beta-\alpha\beta'')\delta'+
(\alpha\beta'-\alpha'\beta)\delta''=\mathfrak{k}
\]
fiat vel $=+1$ vel $=-1$, formamque $f$ transire per substitutionem
\[
\begin{array}{ccc}
\alpha,&\beta,&\delta\\
\alpha',&\beta',&\delta'\\
\alpha'',&\beta'',&\delta''
\end{array}
\]
in
$\left(\begin{smallmatrix}\mathfrak{a},\mathfrak{a}',\mathfrak{a}''\\\mathfrak{b},\mathfrak{b}',\mathfrak{b}''\end{smallmatrix}\right)=\mathfrak{g}$,
cui adiuncta
$\left(\begin{smallmatrix}\mathfrak{A},\mathfrak{A}',\mathfrak{A}''\\\mathfrak{B},\mathfrak{B}',\mathfrak{B}''\end{smallmatrix}\right)=\mathfrak{G}$.
Tunc $g,\mathfrak{g}$ erunt aequivalentes, adeoque etiam $G$ et $\mathfrak{G}$,
et per applicationem praeceptorum in artt.\ 269, 270 traditorum\HistoricalNote{*)}{Eruendo ex transf.\ formae $f$ in $g$, transformationem formae $g$ in $f$; ex hac atque transf.\ formae $f$ in $\mathfrak{g}$, transf.\ formae $g$ in $\mathfrak{g}$; denique ex hac, per transpositionem, transf.\ formae $\mathfrak{G}$ in $G$.}
invenitur, si statuatur
\[
\begin{aligned}
(\beta'\gamma''-\beta''\gamma')\delta+
(\beta''\gamma-\beta\gamma'')\delta'+
(\beta\gamma'-\beta'\gamma)\delta''&=\zeta\\
(\gamma'\alpha''-\gamma''\alpha')\delta+
(\gamma''\alpha-\gamma\alpha'')\delta'+
(\gamma\alpha'-\gamma'\alpha)\delta''&=\eta
\end{aligned}
\]
formam $\mathfrak{G}$ transire in $G$ per substitutionem
\[
\begin{array}{ccc}
k,&0,&0\\
0,&k,&0\\
\zeta,&\eta,&\mathfrak{k}
\end{array}
\]
Hinc erit
\[
B=\eta\mathfrak{k}D+\mathfrak{k}k\mathfrak{B},\qquad
B'=\zeta\mathfrak{k}D+\mathfrak{k}k\mathfrak{B}'
\]
adeoque, propter $\mathfrak{k}k=\pm1$, vel $B\equiv\mathfrak{B}$,
$B'\equiv\mathfrak{B}'$, vel $B\equiv-\mathfrak{B}$, $B'\equiv-\mathfrak{B}'$
$(\mathrm{mod.}\ D)$. In casu priori valores $(B,B')$, $(\mathfrak{B},\mathfrak{B}')$
aequivalentes vocamus, in posteriori oppositos; repraesentationem formae
$\varphi$ autem ad quemlibet valorem
\SourceLeaf{330}
expr.\ $\surd\Delta(p,\allowbreak -q,\allowbreak r)(\mathrm{mod.}\ D)$, qui ex ipsa per methodum in I
deduci potest, \textit{pertinere} dicemus. Hinc omnes valores, ad quos eadem
repraesentatio pertinet, vel aequivalentes erunt vel oppositi.

III.\ Vice versa autem, si ut ante in I repraesentatio formae $\varphi$ per $f$
haec $x=\alpha t+\beta u$ etc.\ ad valorem $(B,B')$ pertinet, qui inde deducitur
adiumento transformationis
\[
\begin{array}{ccc}
\alpha,&\beta,&\gamma\\
\alpha',&\beta',&\gamma'\\
\alpha'',&\beta'',&\gamma''
\end{array}
\]
eadem quoque ad quemvis alium valorem $(\mathfrak{B},\mathfrak{B}')$ pertinebit.
qui illi vel aequivalens est vel oppositus; \textit{i. e.}\ loco ipsorum
$\gamma,\gamma',\gamma''$ alios integros $\delta,\delta',\delta''$ accipere
licebit, pro quibus aequatio $(\Omega)$ haec
\[
(\alpha'\beta''-\alpha''\beta')\delta+
(\alpha''\beta-\alpha\beta'')\delta'+
(\alpha\beta'-\alpha'\beta)\delta''=\pm1
\]
locum habeat, et qui ita comparati sint, ut coëfficiens $4$ et $5$ in forma ei
adiuncta, in quam $f$ per substitutionem $(S)$
\[
\begin{array}{ccc}
\alpha,&\beta,&\delta\\
\alpha',&\beta',&\delta'\\
\alpha'',&\beta'',&\delta''
\end{array}
\]
transit, resp.\ fiant $=\mathfrak{B},\mathfrak{B}'$. Statuatur enim
\[
\pm B=\mathfrak{B}+\eta D,\qquad \pm B'=\mathfrak{B}'+\zeta D
\]
(accipiendo hic et postea signa superiora vel inferiora, prout valores
$(B.\ B')$, $(\mathfrak{B},\mathfrak{B}')$ aequivalentes sunt vel oppositi),
unde $\zeta,\eta$ erunt integri, transeatque $g$ per substitutionem
\[
\begin{array}{ccc}
1,&0,&\zeta\\
0,&1,&\eta\\
0,&0,&\pm1
\end{array}
\]
in formam $\mathfrak{g}$, cuius determinantem esse $\Delta$, in forma adiuncta
vero coëfficientes $4$ et $5$ resp.\ $=\mathfrak{B},\mathfrak{B}'$ fieri facile
perspicietur. Faciendo autem
\[
\alpha\zeta+\beta\eta\pm\gamma=\delta,\quad
\alpha'\zeta+\beta'\eta\pm\gamma'=\delta',\quad
\alpha''\zeta+\beta''\eta\pm\gamma''=\delta''
\]
\SourceLeaf{331}
nullo negotio patebit, $f$ per substitutionem $(S)$ transire in $\mathfrak{g}$,
atque aequationi $(\Omega)$ satisfactum esse. \textit{Q. E. D.}

\Needspace{8\baselineskip}
\Article{283}
Ex his principiis deducitur methodus sequens, omnes repraesentationes proprias
formae binariae
\[
\varphi=ptt+2qtu+ruu
\]
determinantis $D$ per ternariam $f$ determinantis $\Delta$ inveniendi.

I.\ Eruantur omnes valores diversi (\textit{i. e.}\ non-aequivalentes) expressionis
$\surd\Delta(p,\allowbreak -q,\allowbreak r)(\mathrm{mod.}\ D)$. Hoc problema pro eo casu, ubi $\varphi$
est forma primitiva atque $\Delta$ ad $D$ primus, supra (art.\ 233) solutum est,
casusque reliqui ad hunc facillime reducuntur, quam tamen rem fusius hic
explicare brevitas non permittit. Observamus tantummodo, quoties $\Delta$ ad
$D$ primus sit, expressionem $\Delta(p,-q,r)$ residuum quadraticum ipsius $D$
esse non posse, nisi $\varphi$ fuerit forma primitiva. Supponendo enim
\[
\Delta p=BB-DA',\qquad -\Delta q=BB'-DB'',\qquad \Delta r=B'B'-DA
\]
fit
\[
(DB''-\Delta q)^2=(DA'+\Delta p)(DA+\Delta r)
\]
hinc, per evolutionem et substituendo $qq-pr$ pro $D$, fit
\[
(qq-pr)(B''B''-AA')-\Delta(Ap+2B''q+A'r)+\Delta\Delta=0
\]
unde facile concluditur, si $p,q,r$ divisorem communem haberent, hunc etiam
ipsum $\Delta\Delta$ metiri; tunc vero $\Delta$ ad $D$ primus esse non posset.
Quare $p,q,r$ divisorem communem habere nequeunt, sive $\varphi$ erit forma
primitiva.

II.\ Designemus multitudinem horum valorum per $m$, supponamusque.
inter eos reperiri $n$ valores, qui sibi ipsis oppositi sint (statuendo $n=0$,
quando tales non adsunt). Tunc manifestum est, ex $m-n$ reliquis valoribus binos
semper oppositos fore (quoniam cuncti valores complete haberi supponuntur);
reiiciatur e binis quibusque valoribus oppositis unus ad libitum, remanebuntque
omnino valores $\frac{1}{2}(m+n)$. Ita \textit{e. g.}\ ex octo valoribus expr.\
$\surd-1(19,-3,41)$
